% ─────────────────────────────────────────────────────────────
%  01-minimal.tex — un premier document LaTeX complet
%  À compiler avec pdfLaTeX (réglage par défaut d'Overleaf).
% ─────────────────────────────────────────────────────────────
\documentclass[12pt]{article}   % classe du document et taille du texte

% ── Préambule : réglages et extensions (packages) ────────────
\usepackage[T1]{fontenc}         % accents et césures corrects
\usepackage[french]{babel}       % typographie française
\usepackage{amsmath, amssymb}    % outils mathématiques
\usepackage{siunitx}             % nombres et unités
\sisetup{output-decimal-marker = {,}}   % virgule décimale
\usepackage[margin=2.5cm]{geometry}     % marges de la page

\title{Mon premier document}
\author{Prénom Nom}
\date{\today}

% ── Corps du document ────────────────────────────────────────
\begin{document}

\maketitle

\section{Du texte}

Un paragraphe se termine par une ligne vide.
Un simple retour à la ligne, comme ici, ne change rien.

Voici un nouveau paragraphe. Le symbole pour cent s'écrit \%,
car le symbole seul commence un commentaire. % Ceci est un commentaire.

\section{Des mathématiques}

Dans le texte, on écrit $a^2 + b^2 = c^2$ entre deux symboles dollar.
Une équation importante se met en évidence :
\[
  x = \frac{-b \pm \sqrt{b^2 - 4ac}}{2a}.
\]
Un nombre décimal : $\num{3,75}$ ; une mesure : $\qty{12,5}{\centi\metre}$.

\section{Une liste}

\begin{enumerate}
  \item Lire l'énoncé.
  \item Faire un schéma.
  \item Résoudre et valider.
\end{enumerate}

\end{document}
